\section{Related work}
\label{sec:related}

We organise prior work along two axes: (i)~whether the model is asked to \emph{interpret} a diagram or \emph{produce} one, and (ii)~whether grading is symbolic, vision--language, or fixed-answer.
\Cref{tab:benchmark_comparison} positions \GeoGenBench against the closest prior benchmarks.

\begin{table}[t]
\caption{\GeoGenBench vs.\ closest prior benchmarks. \emph{Auto-grade}: grading does not require human or MLLM judgment.}
\label{tab:benchmark_comparison}
\centering
\small
\begin{tabular}{lllrlll}
\toprule
\textbf{Benchmark} & \textbf{Year} & \textbf{Task} & \textbf{\# items} & \textbf{Auto-grade} & \textbf{Modality} \\
\midrule
Geometry3K \citep{lu2021intergps} & 2021 & QA & 3{,}002 & yes (MC) & image+text \\
GeoQA \citep{chen2021geoqa} & 2021 & QA (program) & 4{,}998 & yes (exec) & image+text \\
UniGeo \citep{chen2022unigeo} & 2022 & QA + proving & 14{,}541 & yes (seq) & image+text \\
MathVista \citep{lu2024mathvista} & 2024 & QA (mixed) & 6{,}141 & partial & image+text \\
GeoEval \citep{zhang2024geoeval} & 2024 & QA + variants & 5{,}050 & yes & image+text \\
GeoBench \citep{geobench2026} & 2026 & hierarchical QA & 1{,}021 & yes (verified) & image+text \\
T2I-CompBench++ \citep{huang2025t2icompbenchpp} & 2025 & T2I gen & 8{,}000 & partial (VLM) & text $\to$ image \\
GenExam \citep{wang2025genexam} & 2025 & T2I exam gen & 1{,}000 & LLM-judge & text $\to$ image \\
\midrule
\textbf{\GeoGenBench (ours)} & \textbf{2026} & \textbf{diagram gen} & \textbf{801} & \textbf{yes (symbolic)} & \textbf{text $\to$ TikZ $\to$ SVG} \\
\bottomrule
\end{tabular}
\end{table}

\paragraph{Geometry QA benchmarks.}
Geometry3K \citep{lu2021intergps}, GeoQA \citep{chen2021geoqa}, UniGeo \citep{chen2022unigeo}, MathVista \citep{lu2024mathvista}, GeoEval \citep{zhang2024geoeval}, and GeoBench \citep{geobench2026} all measure whether a model can \emph{interpret} a given diagram and answer a question about it.
\GeoGenBench instead measures whether a model can \emph{produce} a correct diagram from a textual description --- the deployed task when LLMs author worksheets, slides, and technical illustrations.

\paragraph{Generation-side benchmarks.}
T2I-CompBench++ \citep{huang2025t2icompbenchpp} and GenExam \citep{wang2025genexam} grade generated images via vision--language judges (CLIP, BLIP, MLLMs) or LLM judges; these cannot tell whether a ``right triangle'' actually has a right angle, only whether the image \emph{looks like} one.
\GeoGenBench grades in symbolic space (TikZ $\to$ compiled SymPy geometry) where exact verification is possible: a \propname{right\_angle\_at\_C} verdict is \texttt{Triangle(B, C, A).angles[C] == $\pi/2 \pm \tau$}, not a judge's opinion.
GenExam is the closest prior art (generative-exam framing); we differ on grading mechanism (deterministic SymPy vs MLLM judge, $\sim$$1000\times$ cheaper and fully reproducible) and modality (TikZ vs raster).

\paragraph{Symbolic / formal grading.}
Lean and Isabelle corpora (miniF2F, ProofNet) and AlphaGeometry's neuro-symbolic deduction engine \citep{trinh2024alphageometry} apply machine-checked grading to \emph{proofs}; \GeoGenBench extends the automatic-verification tradition into \emph{diagram generation}, where the verifiability target is a geometric configuration rather than a propositional proof.
\GeoGenBench occupies the previously-empty cell of \emph{generation-side, automatically-graded} geometry, enabled by the choice to grade in symbolic space rather than pixel space.
